% !TeX root = ../induk.tex
\begin{table}[H]
\centering
\caption{Matriks Transisi IPM Tahun 2020 ke 2021} \label{tab:ipm_transisi_2020_2021}
\begin{tabular}{|c|p{2cm}|p{2cm}|p{2cm}|p{2cm}|}
\hline
\makecell{\textbf{Tahun 2020} \\ \textbf{2021}} & \textbf{Sangat Tinggi (2021)} & \textbf{Tinggi (2021)} & \textbf{Sedang (2021)} & \textbf{Rendah (2021)} \\ \hline
\textbf{Sangat Tinggi (2020)} & Kota Jakarta Selatan, Kota Jakarta Timur, Kota Jakarta Pusat, Kota Jakarta Barat, Kota Jakarta Utara, Kota Bandung, Kota Bekasi, Kota Depok, Kota Tangerang Selatan & - & - & - \\ \hline
\textbf{Tinggi (2020)} & - & Kep. Seribu, Bogor, Bandung, Ciamis, Sumedang, Purwakarta, Karawang, Bekasi, Kota Bogor, Kota Sukabumi, Kota Cirebon, Kota Cimahi, Kota Tasikmalaya, Kota Banjar, Tangerang, Kota Tangerang, Kota Cilegon, Kota Serang & - & - \\ \hline
\textbf{Sedang (2020)} & - & - & Sukabumi, Cianjur, Garut, Tasikmalaya, Kuningan, Cirebon, Majalengka, Indramayu, Subang, Bandung Barat, Pangandaran, Pandeglang, Lebak, Serang & - \\ \hline
\textbf{Rendah (2020)} & - & - & - & - \\ \hline
\end{tabular}
\end{table}